% ===== 底本 第 2 章《拉格朗日环面纤维化》；片段 tr02a；源块 B038--B085 =====

\section{拉格朗日环面纤维化}

% [源 B038]
在本讲中，我们将回顾关于拉格朗日环面纤维化（Lagrangian torus fibration）的一些基本几何事实，特别是正则的情形。

% [源 B039]
\subsection{正则拉格朗日纤维化}

% [源 B040]
本讲的目标是建立正则拉格朗日纤维化的下列性质：

\begin{itemize}[leftmargin=1.4em,itemsep=1pt,topsep=2pt]
  % [源 B041]
  \item 它们在局部上以完全可积哈密顿系统（completely integrable hamiltonian system）为模型；
  % [源 B042]
  \item 正则拉格朗日纤维化的纤维是拉格朗日环面；
  % [源 B043]
  \item 正则拉格朗日纤维化的底空间带有整仿射结构（integral affine structure）。
\end{itemize}

% [Tutor Note: B043 源文 “Basis of regular Lagrangian fibrations” 中的 Basis 指纤维化的底空间（base），不是线性代数意义的基；故译作“底空间”。]

% [源 B044]
我们先从拉格朗日纤维化的定义说起。

% [源 B045]
\begin{definition*}
拓扑空间 $B$ 的一个\emph{分层}（stratification）是一个滤列
\[
  \emptyset =B_{-1}\subseteq B_0\subseteq B_1\subseteq\dotsb\subseteq B_{d+1}\subseteq\dotsb\subseteq B
\]
其中每个 $B_d$ 都是闭子集，使得 $\forall d$，$S_d (B):= B_d\setminus B_{d-1}$ 是 $d$ 维光滑流形，且 $\displaystyle B=\bigcup_{d\geq 0} B_d$。若当 $d\gg 0$ 时 $S_d(B)=\emptyset$，则称 $B$ 是\emph{有限维的}；若 $B$ 有限维，且 $n$ 是使 $S_n (B)\neq\emptyset$ 的最大整数，则称 $B$ 是\emph{$n$ 维的}。此时称 $S_n(B)$ 为\emph{顶层}（top stratum）。
\end{definition*}

% [Tutor Note: B045 的滤列自 $B_{-1}=\emptyset$ 起步，此步不可省：否则 $d=0$ 时 $S_0(B)=B_0\setminus B_{-1}$ 无定义。核对 $d=0$ 的情形：$S_0(B)$ 应为若干孤立点，恰是 0 维光滑流形。]

% [源 B046]
现在我们准备给出拉格朗日纤维化的定义。

% [源 B047]
\begin{definition*}
一个\emph{拉格朗日纤维化}（Lagrangian fibration）是映射 $(M^{2n},\omega)\xrightarrow{\pi} B^n$，使得

\begin{itemize}[leftmargin=1.4em,itemsep=1pt,topsep=2pt]
  % [源 B048]
  \item $(M^{2n},\omega)$ 是辛流形（symplectic manifold），且 $B$ 是 $n$ 维分层空间；
  % [源 B049]
  \item $\pi$ 固有（proper）且连续；
  % [源 B050]
  \item 在顶层之上，$\pi$ 是光滑淹没（submersion），其纤维为连通的拉格朗日纤维。其余纤维连通，其各分层均为迷向（isotropic）的。
\end{itemize}

% [源 B051]
% [Tutor Note: B051 在底本 HTML 中位于同一个 Definition 的 div 内（紧接上述 ul 之后），故并入本定义环境，不另起段落。]
记 $B^{reg}:=S_n (B)$ 为顶层，亦即 $\pi$ 的\emph{正则轨迹}（regular locus）；记 $B^{sing}=B\setminus B^{reg}$ 为 $\pi$ 的\emph{判别轨迹}（discriminant locus），余维数 $\geq 1$。所谓\emph{正则/光滑}拉格朗日纤维化，是指 $B=B^{reg}$，即不存在更低维的分层。
\end{definition*}

% [源 B052]
\subsection{可积哈密顿系统}

% [源 B053]
历史上，拉格朗日纤维化源于对经典哈密顿系统动力学性质的研究。哈密顿量是定义在相空间上的光滑函数，充当运动系统的不变量；哈密顿函数的存在使我们能够“降低系统的自由度”，意即质点必须沿着由哈密顿函数决定的等值超曲面运动。本小节我们将描述这样的哈密顿函数组：它们允许我们归纳地施行降维，并说明它们与拉格朗日纤维化的联系。

% [Tutor Note: B053 的 “level hypersurfaces” 指等值超曲面 $H^{-1}(c)$，余维为 1，因为单个哈密顿量只提供一个守恒量；这正说明为何要用多个 $H_i$ 才能把自由度逐个降下去。]

% [源 B054]
\begin{definition*}
设 $(M^{2n},\omega)$ 是辛流形。函数 $\vec{H}=(H_1,\dotsb, H_k)\colon M\to\mathbb{R}^k$ 称为\emph{可积哈密顿系统}（integrable hamiltonian system），若

\begin{itemize}[leftmargin=1.4em,itemsep=1pt,topsep=2pt]
  % [源 B055]
  \item 各 $H_i$ 两两\emph{Poisson 交换}，即对一切 $i,j$ 有 $\{H_i, H_j\}=0$。（其中 $\{H_i,H_j\}=\omega (X_{H_i}, X_{H_j})$）
  % [源 B056]
  \item $\vec{H}(X)$ 在其闭包中包含正则值的一个开稠子集。（即 $dH_1,\dotsb, dH_k$ 线性无关）
\end{itemize}

% [源 B057]
若 $k=n$，且进一步满足

\begin{itemize}[leftmargin=1.4em,itemsep=1pt,topsep=2pt]
  % [源 B058]
  \item $\vec{H}$ 固有，且纤维连通。
\end{itemize}

则称 $\vec{H}$ 是\emph{完全可积的}（completely integrable）。
\end{definition*}

% [Tutor Note: B056 中的 $X$ 未在文中定义，按上下文应为 $M$（或 $\vec{H}(M)$）；“in its closure” 指 $\vec{H}(M)$ 的闭包。照译，不改符号。]

% [源 B059]
\begin{proposition*}
哈密顿函数 $H$ 的哈密顿流保持在 $H$ 的水平集 $H^{-1} (c)$ 之内。
\end{proposition*}

% [源 B060]
\begin{proof}
设 $\gamma(t)$ 是 $H$ 的积分曲线，则 $\displaystyle\frac{d}{dt} H(\gamma(t))=dH(\dot{\gamma} (t))=dH(X_H(\gamma(t)))=0$。
\end{proof}

% [Tutor Note: B060 源文 “the integral curve of $H$” 实指哈密顿向量场 $X_H$ 的积分曲线（$H$ 是函数，谈不上积分曲线）；末式 $dH(X_H)=0$ 也印证了这一点。照译。]

% [源 B061]
\begin{proposition*}
若 $H,G$ 是两个 Poisson 交换的哈密顿函数，则它们的哈密顿流交换。
\end{proposition*}

% [源 B062]
\begin{proof}
$\left[X_{H_i},X_{H_j}\right]=X_{\{H_i,H_j\}}=0$。
\end{proof}

% [Tutor Note: B062 的命题用 $H,G$，证明却写成下标 $H_i,H_j$；应读作 $H=H_i$、$G=H_j$，左端为向量场的交换子。照译，不改符号。]

% [源 B063]
由 Frobenius 定理立即可知

% [源 B064]
\begin{corollary*}
完全可积哈密顿系统 $\vec{H}\colon M\to\mathbb{R}^n$ 在 $M$ 上生成一个 $\mathbb{R}^n$-作用。
\end{corollary*}

% [源 B065]
{\sloppy 设 $c$ 是 $\vec{H}$ 的正则值，则由完全可积哈密顿系统的定义立即可知 $TH^{-1} (c)=\operatorname{span}\left\{X_{H_1},\dotsb,X_{H_n}\right\}$，且 $H^{-1} (c)$ 是 $\mathbb{R}^n$-作用的一条拉格朗日轨道。现在取 $\pi$ 的一个局部截面 $\sigma\colon B\to M$，其中 $B\subseteq\vec{H}(M)\subseteq\mathbb{R}^n$ 是正则值的开子集。\par}

% [Tutor Note: 底本源文用页面自定义宏 `\espan`（HTML 中 tex.macros 定义为 \operatorname{span}），并非笔误；本译本按 \operatorname{span} 排印，数学含义不变。]
% [Tutor Note: B065 的 $H$ 指映射 $\vec{H}$（与 B059 中标量函数 $H$ 的记号混用），$TH^{-1}(c)$ 读作水平集 $H^{-1}(c)$ 的切空间。照译。]

% [源 B066]
\begin{definition*}
给定完全可积哈密顿系统 $\vec{H}\colon M\to\mathbb{R}^n$ 与局部截面 $\sigma\colon B\to M$，在 $b\in B$ 处的\emph{周期格}（period lattice）定义为如下子集：
\[
  \Lambda_b^H :=\{t\in\mathbb{R}^n\vert \phi_H^t (\sigma (b))=\sigma (b)\}.
\]
我们另外把\emph{周期格}定义为 $\mathbb{Z}$-局部系统 $\Lambda^H =\{(b,t)\in B\times\mathbb{R}^n\vert t\in\Lambda_b^H\}$。若 $\Lambda =B\times (2\pi\mathbb{Z})^n$，则称它为\emph{标准的}。
\end{definition*}

% [Tutor Note: B066 的两个 $\Lambda$ 层次不同：$\Lambda_b^H$ 是单点的格，$\Lambda^H$ 是拼起来的 $\mathbb{Z}$-局部系统；“standard” 指 $\Lambda=B\times(2\pi\mathbb{Z})^n$，即 $\phi_H^t$ 的周期恰为 $2\pi$ 的整格，这个 $2\pi$ 约定在 B072 的 $2\pi A(b)^T\mathbb{Z}^n$ 与 B103 的通量计算中再次出现。]

% [源 B067]
\begin{theorem*}[Little Arnold--Liouville 定理]
设 $\vec{H}\colon M\to\mathbb{R}^n$ 是 $M$ 上的可积哈密顿系统。设 $\Omega (\sigma (b))$ 是 $\sigma (b)$ 在 $\mathbb{R}^n$-作用下的轨道，则对某个 $k$，$\Omega (\sigma (b))$ 微分同胚于 $\displaystyle\left(\mathbb{R}^k/\mathbb{Z}^k\right)\times\mathbb{R}^{n-k}$；并且当 $\Omega (\sigma (b))$ 紧致时，它是一个环面。
\end{theorem*}

% [Tutor Note: B067 的证明被原文略去，可补出如下：由 B064 的 $\mathbb{R}^n$-作用与 Frobenius 定理，轨道 $\Omega(\sigma(b))$ 同构于 $\mathbb{R}^n/\Lambda_b^H$；$\Lambda_b^H$ 是秩 $k$ 的格，故同构于 $(\mathbb{R}^k/\mathbb{Z}^k)\times\mathbb{R}^{n-k}$，且紧致当且仅当 $k=n$（此时得 $n$ 环面）。]

% [源 B068]
\begin{exercise*}
Little Arnold--Liouville 定理假定了可积系统 $\vec{H}$ 存在局部截面 $\sigma\colon B\to M$，但实际上可以证明这样的截面在局部总是存在的。例如，若 $x\in M$ 是 $\vec{H}$ 的正则点，则总存在 $x$ 的开邻域 $U$，使 $\vec{H}(U)=B$ 在 $\mathbb{R}^n$ 中为开集，且存在 $\vec{H}$ 的局部拉格朗日截面 $\sigma\colon B\to U$ 满足 $\sigma (\vec{H} (x))=x$。
\end{exercise*}

% [Tutor Note: B068 的构造即取 B069 的 Liouville 映射 $\psi\colon B\times\mathbb{R}^n\to M$ 并令 $\sigma:=\psi(\cdot,0)$：$\psi^{\ast}\omega=\sum db_i\wedge dt_i$ 在 $t=0$ 的截面上消失，故 $\sigma(B)$ 是拉格朗日子的，再经平移即可满足 $\sigma(\vec{H}(x))=x$。]

% [源 B069]
在拉格朗日截面 $\sigma\colon B\to M$ 的局部附近，我们可以用 $\psi (b,t)=\phi_H^t (\sigma(b))$ 构造映射 $\psi\colon B\times\mathbb{R}^n\to M$；可以验证它既是浸入又是淹没，且满足 $\psi^{\ast}\omega =\displaystyle\sum_{i=1}^n db_i\wedge dt_i$。这个映射称为\emph{Liouville 坐标}（Liouville coordinate）。

% [Tutor Note: B069 的 $\psi^{\ast}\omega=\sum_{i=1}^n db_i\wedge dt_i$ 可按两个方向核对：沿 $t$ 方向 $\psi$ 是哈密顿流的积分，保 $\omega$；沿 $b$ 方向，$dt_i$ 与 $\partial_{b_i}$ 的配对给出 $\omega(\partial_{b_i},\partial_{t_j})=\delta_{ij}$。取 $n=1$ 即退化为面积形式 $db\wedge dt$。]

% [源 B070]
\begin{remark*}
若 $\sigma$ 只是局部截面，则上述构造仍是局部微分同胚，但 $\psi^{\ast}\omega =\sum db_i\wedge dt_i +\beta$，其中 $\beta$ 是某个闭 $2$-形式。只有当 $\sigma$ 是拉格朗日的，才有 $\beta=0$。
\end{remark*}

% [Tutor Note: B070 中 $\beta$ 确实闭：$d\beta=d(\psi^{\ast}\omega)=\psi^{\ast}d\omega=0$；故原文称其为闭 $2$-形式是自洽的。]

% [源 B071]
\begin{theorem*}[作用-角坐标]
给定 $\vec{H}\colon M\to B\subseteq\mathbb{R}^n$，其中 $B$ 是 $n$ 维圆盘且其上的纤维都是正则纤维，并给定拉格朗日截面 $\sigma\colon B\to M$，则存在局部坐标变换 $\alpha\colon B\to C\subseteq\mathbb{R}^n$，使得对 $a:=\alpha\circ\vec{H}\colon M\to C\subseteq\mathbb{R}^n$ 而言，周期格是标准的。
\end{theorem*}

% [源 B072]
\begin{proof}
由于 $B$ 可缩，我们可以取 $\Lambda^H$ 在 $B$ 上的一组基，其周期为 $2\pi\tau_1 (b),\dotsb, 2\pi\tau_n (b)$，使得 $\Lambda_b^H =2\pi A(b)^T\mathbb{Z}^n$，其中 $A(b)$ 是以 $\tau_i (b)$ 为元素的对角矩阵。现在我们寻找满足偏微分关系 $\displaystyle\frac{\partial\alpha_i}{\partial b_j}=A_{ij}$ 的映射 $\alpha$；该关系可解当且仅当 $\displaystyle\frac{\partial A_{ij}}{\partial b_k}=\frac{\partial A_{ik}}{\partial b_j}$。而这一点可以用 Liouville 坐标验证。
\end{proof}

% [Tutor Note: B072 的可解性条件即混合偏导可交换（Clairaut）：$\partial_{b_k}\partial_{b_j}\alpha_i$ 的对称性给出 $\partial_{b_k}A_{ij}=\partial_{b_j}A_{ik}$。作者称“可用 Liouville 坐标验证”，指在 B069 的坐标里辛形式已是标准形，$\Lambda_b^H$ 的基与 $\partial_{t_i}$ 的关系直接给出 $A$。核对 $n=1$：$\alpha'(b)=\tau(b)$ 总有解，标准格为 $2\pi\mathbb{Z}$。]

% [源 B073]
我们称 $a$ 为\emph{作用坐标}（action coordinate），$t$ 为\emph{角坐标}（angle coordinate）。

% [Tutor Note: B073 的 $t$ 取值为 $\mathbb{R}^n/(2\pi\mathbb{Z})^n=\mathbb{T}^n$（由 B066 的 standard 周期格），这正是 B074 中局部模型 $B\times\mathbb{T}^n$ 的第二个因子。]

% [源 B074]
\begin{corollary*}[Big Arnold--Liouville 定理]
设 $\vec{H}\colon M\to\mathbb{R}^n$ 完全可积，则所有正则纤维都是拉格朗日环面，且 $M$ 局部辛同胚于 $B\times\mathbb{T}^n$。
\end{corollary*}

% [源 B075]
现在我们可以证明拉格朗日纤维化的第一条重要性质：

% [源 B076]
\begin{proposition*}
正则拉格朗日纤维化在局部上以完全可积哈密顿系统为模型。
\end{proposition*}

% [源 B077]
\begin{proof}
给定正则值的一个开子集及其坐标卡 $\phi\colon U_b\to\mathbb{R}^n$，我们可以作复合映射 $\pi^{-1} (U_b)\xrightarrow{\phi\circ\pi}\mathbb{R}^n$，可以验证它是完全可积哈密顿系统。
\end{proof}

% [Tutor Note: B077 的“可验证”具体为：取 $H_i:=\phi_i\circ\pi$，则各 $dH_i$ 沿纤维方向为零（纤维落在 $\pi$ 的原像中），而在淹没处 $dH_1,\dotsb,dH_n$ 线性无关；纤维连通性由 B047 定义中拉格朗日纤维化的第三条（connected fibers）给出。]

% [源 B078]
\begin{corollary*}
正则拉格朗日纤维化的纤维是环面。
\end{corollary*}

% [源 B079]
作用-角坐标是局部的，但我们可以把它们拼接起来：这样就能把可缩空间上的正则拉格朗日纤维化粘成某个光滑流形 $B$，它带有一种特定的刚性结构，我们称这样的流形为\emph{整仿射流形}（integral affine manifold）。

% [源 B080]
\begin{definition*}
一个\emph{整仿射变换}（integral affine transformation）是映射 $T\colon\mathbb{R}^n\to\mathbb{R}^n$，满足 $T(b)=Ab+c$，其中 $A\in GL(n,\mathbb{Z})$、$c\in\mathbb{R}^n$。$B$ 上的一个\emph{整仿射结构}（integral affine structure）由一张图册给出，其转移函数都是整仿射变换。我们把容许整仿射结构的流形称为\emph{整仿射流形}。
\end{definition*}

% [Tutor Note: B080 的整仿射变换要求线性部分落在 $GL(n,\mathbb{Z})$（保标准格 $\mathbb{Z}^n$ 与 $2\pi$ 周期），这与 B072 中 $\Lambda_b^H=2\pi A(b)^T\mathbb{Z}^n$ 的取法一致；核对：$\det A=\pm1$，否则格会被放大或缩小。]

% [源 B081]
\begin{lemma*}
设 $G,H\colon M\to\mathbb{R}^n$ 为完全可积哈密顿系统，$\psi\colon\mathbb{R}^n$ 为微分同胚，且使下面的交换图交换：

% [源 B082]
\begin{tikzcd}
  & M \arrow[dl, "G"'] \arrow[dr, "H"] & \\
  \R^n \arrow[rr, "\psi"'] & & \R^n
\end{tikzcd}

% [源 B083]
则 $\psi$ 是整仿射的。
\end{lemma*}

% [Tutor Note: B081 源文作 $\psi\colon\mathbb{R}^n$，缺靶空间；按交换图 B082 应读作 $\psi\colon\mathbb{R}^n\to\mathbb{R}^n$。照译，不改原文，详见回执。]

% [源 B084]
\begin{proof}
由于必须有 $d_b\psi\colon\Lambda_{\psi (b)}^G\to\Lambda_b^H$，对每个 $b$，$d_b\psi$ 形如 $d_b\psi (\xi)=A(b)\xi +c(b)$，其中矩阵 $A(b)$ 落在 $GL(n,\mathbb{Z})$ 中，且 $c(b)\in\mathbb{R}^n$。由于 $G$ 与 $H$ 的周期格都是标准的，$d_b\psi$ 必为常值，从而 $\psi (b)=Ab+c$。
\end{proof}

% [Tutor Note: B084 的论证用到“周期格是标准的”这一陈述之外的额外假设（引理陈述中并没有）；实际做法是先用 B071 的作用-角坐标把 $G$ 与 $H$ 两侧都化为标准格，才能断言 $d_b\psi$ 为常值。另注：$d_b\psi (\xi)=A(b)\xi +c(b)$ 严格说 $c(b)$ 应并入 $\psi$ 的平移项（导数 $d_b\psi$ 本身是线性的），原文照译。]

% [源 B085]
\begin{corollary*}
若 $\pi\colon M\to B$ 是正则拉格朗日纤维化，则 $B$ 继承一个整仿射结构。
\end{corollary*}

% [Tutor Note: B085 的证明路线：B076（局部以可积系统为模型）+ B071（作用-角坐标）+ B081（两套标准周期格之间的过渡映射是整仿射的）三者拼接，得到 $B$ 上的整仿射图册；核对过渡函数：其形如 $Ab+c$，$A\in GL(n,\mathbb{Z})$，与 B080 的定义一致。]
